\chapter{\ifenglish Introduction\else บทนำ\fi}

\section{\ifenglish Project rationale\else ที่มาของโครงงาน\fi}
ในปัจจุบัน การโหวตนั้นเป็นสิ่งที่พบเห็นได้ทั่วไปเนื่องจากว่าประเทศไทยนั้นมีการปกครองแบบประชาธิปไตร แต่ว่าปัญหาในการโหวตนั้นมีให้เห็นอยู่มากมายเช่น บัตรเสีย การนับคะแนนที่ผิดพลาด
แต่หากเรามาดูรูปแบบการแก้ปัญหาก็พบว่ามีการนำเครื่องเลือกตั้งแบบอิเล็กโทรนิคมาใช้แต่ว่าปัญหาของมันก็คือการที่มีความเสี่ยงจะถูกโจมตีจาก Man-in-the-Middle (MITM) พอมาดูรูปแบบการโหวตต่อไปก็คือการลงคะแนนเสียงออนไลน์ ซึ่งมีความปลอดภัยสูงเพราะว่ามีการเข้ารหัสแบบ RSA ซึ่งเป็นการเข้ารหัสแบบ Asymmetric Encryption แต่ว่าการมาถึงของ
ควอนตัมคอมพิวเตอร์นั้นจะทำให้ การเข้ารหัสแบบนี้ไม่ปลอดภัยเนื่องจากว่าควอนตัมคอมพิวเตอร์นั้นสามารถรัน Shor's algorithm ได้ทำให้การโหวตแบบออนไลน์นั้นไม่ปลอดภัย โครงงานวิจัยเราจึงจัดทำขึ้นเพื่อนำเสนอการส่ง key โดยใช้งาน Quantum Key Distribution (QKD) ในการส่ง key โดยเราจะทำการเข้ารหัสผลการเลือกตั้งเป็นแบบ Symmetric Encryption จากนั้นเราก็จะทำการส่ง key ผ่านทาง Quantum Key Distribution (QKD) ซึ่งรูปแบบนี้เป็นวิธีที่ปลอดภัยต่อการมาถึงของเครื่องควอนตัมคอมพิวเตอร์เพราะว่าคุณสมบัติของควอนตัมนั้นไม่สามารถทำซ้ำได้หากมีการดักฟังก็จะทำให้ผลลัพธ์เปลี่ยน

\section{\ifenglish Objectives\else วัตถุประสงค์ของโครงงาน\fi}
\begin{enumerate}
    \item เพื่อประยุกต์การใช้งาน Quantum Key Distribution ในการส่ง key
    \item เพื่อจำลองการโหวตโดยใช้ควอนตัมคอมพิวเตอร์โดยมีสถานการคือส่งได้แบบปกติและกำลัง
ถูกดักฟัง
\end{enumerate}

\section{\ifenglish Project scope\else ขอบเขตของโครงงาน\fi}

\subsection{\ifenglish Hardware scope\else ขอบเขตด้านฮาร์ดแวร์\fi}
เนื่องจากว่าควอนตัมคอมพิวเตอร์นั้นยังไม่ได้มีการใช้งานอย่างแพร่หลาย เราจึงใช้เครื่องคอมพิวเตอร์ทั่วไปที่สามารถต่ออินเทอร์เน็ตได้ในการเข้าถึงเครื่องควอนตัมคอมพิวเตอร์

\subsection{\ifenglish Software scope\else ขอบเขตด้านซอฟต์แวร์\fi}
จากข้อที่ผ่านมาได้กล่าวว่ายังไม่มีเครื่องควอนตัมคอมพิวเตอร์ในการใช้งาน เราจึงต้องมีการใช้งานแพลตฟอร์ม Quantum Simulation ในออนไลน์ เช่น IBM Quantum และได้มีการใช้งาน ibrary qiskit\_ibm\_runtime ของ Python เพื่อเขียนโปรแกรม Quantum Assembly 

\section{\ifenglish Expected outcomes\else ประโยชน์ที่ได้รับ\fi}
\begin{enumerate}
    \item  ได้ใช้ความรู้เกี่ยวกับการเข้ารห้สแบบ Symmetric Encryption 
    \item ได้จำลองการส่ง key โดยใช้ Quantum Key Distribution (QKD)
    \item  ได้ทำแบบจำลองรูปแบบการเลือกตั้งโดยใช้ควอนตัมคอมพิวเตอร์
\end{enumerate}

\section{\ifenglish Technology and tools\else เทคโนโลยีและเครื่องมือที่ใช้\fi}

\subsection{\ifenglish Hardware technology\else เทคโนโลยีด้านฮาร์ดแวร์\fi}
เครื่องคอมพิวเตอร์ส่วนบุคคลที่มีความสามารถในการเชื่อมต่ออินเทอร์เน็ต

\subsection{\ifenglish Software technology\else เทคโนโลยีด้านซอฟต์แวร์\fi}
\begin{enumerate}
    \item ภาษา phyton และ library qiskit\_ibm\_runtime 
    \item database mongodb
    \item ภาษา typescript และใช้งาน library ReactJS
    \item web browser สำหรับเข้าใช้งาน IBM Quantum
\end{enumerate}

\section{แผนการะใช้งบประมาณ}
เนื่องจากว่าโครงงานนี้เป็นการทำงานเกี่ยวกับการใช้โปรแกรมจึงเลยไม่จำเป็นต้องใช้งบประมาณแต่ว่าอาจจะมีการใช้งบประมาณในการจ่ายค่าใช้บริการแพลตฟอร์ม IBM Quantum

\section{\ifenglish Project plan\else แผนการดำเนินงาน\fi}
\begin{plan}{6}{2020}{2}{2021}
    \planitem{7}{2020}{8}{2020}{ศึกษาค้นคว้า}
    \planitem{8}{2020}{1}{2021}{ชิล}
    \planitem{2}{2021}{2}{2021}{เผา}
    \planitem{12}{2019}{1}{2022}{ทดสอบ}
\end{plan}

\section{\ifenglish Roles and responsibilities\else บทบาทและความรับผิดชอบ\fi}
\begin{enumerate}
    \item นายภคิน นิรมิตสีมา หน้าที่ความรับผิดชอบคือเป็นคนทำระบบของ Tally Server ซึ่งเป็นระบบในการจัดเก็บผลการโหวตของคนที่มาโหวตไว้โดยที่ไม่สามารถระบุตัวคนโหวตได้
    \item นายยุทธกานต์ สาจุ้ย หน้าที่ความรับผิดชอบคือทำระบบหน้าบ้านซึ่งเป็นเว็บในการทดสอบระบบการโหวตโดยใช้งานควอนตัมคอมพิวเตอร์และเป็นคนทำระบบ Authentication Server ซึ่งเป็นระบบในการ
ยืนยันตัวตนสำหรับผู้มีสิทธ์ในการโหวต
    \item อาทินันท์ แจ้งสว่าง หน้าที่ความรับผิดชอบคือการทำ Simulator เพื่อแสดงให้เห็นว่าการเข้ารหัสของควอนตัมคอมพิวเตอร์นั้นสามารถทำงานได้ถูกต้อง
    \item ไพสิฐ เลิศอนันต์พิพัฒน์ หน้าที่ความรับผิดชอบคือการออกแบบหน้าตาของเว็บแอปที่ใช้ในการ
จำลองการโหวตโดยใช้งานควอนตัมคอมพิวเตอร์ซึ่งจะมีทั้งแบบโดนดักฟังและไม่ดักฟัง รวมถึงการ
ออกแบบและใช้งาน Quantum Key Distribution ในการส่ง key
\end{enumerate}

\section{\ifenglish%
Impacts of this project on society, health, safety, legal, and cultural issues
\else%
ผลกระทบด้านสังคม สุขภาพ ความปลอดภัย กฎหมาย และวัฒนธรรม
\fi}

โครงงานนี้มีผลกระทบโดยตรงเกี่ยวกับการศึกษาเพราะว่าเป็นการนำองค์ความรู้ที่เกี่ยวกับด้านควอนตัมมาประยุกต์ใช้งานในการที่จะทำให้ระบบโหวตมีความปลอดภัยมากขึ้น ต่อมาคือในด้านความปลอดภัยจะทำให้การโจมตีที่ใช้งานควอนตัมคอมพิวเตอร์ไม่สามารถโจมตีได้และสุดท้ายคือจะทำให้หน่วยงานที่เกี่ยวกับการเลือกตั้งสามารถนำระบบการโหวตของเรานั้นไปประยุกต์ใช้เพื่อให้มีความโปร่งใสและสามารถตรวจสอบผลการลงคะแนนของบุคคลได้




